%% Appendix A: Mathematical Derivations: appendices/appendixA.tex

% ==============================================================================
% CONTENT BLUEPRINT — Appendix A: Mathematical Derivations
% ==============================================================================
% Mathematical derivations belonging to the main text. Reader can skip
% Appendix A on first read; should be self-contained on second read.
%
% ----------------------------------------------------------------------------
% Recommended content:
%   • Full step-by-step derivations of every non-trivial identity used in
%     Chapters 2–3 (Hohenberg-Kohn, Kohn-Sham, equation-of-state fit,
%     dielectric function from interband transitions, etc.).
%   • Each derivation:
%       – Clearly state starting point and assumption list at the top.
%       – Every step uses a label like \stepref{sym:1} or simply unnumbered
%         aligned equations. Number only the final result.
%       – End with a boxed "Result" equation that is the working formula
%         used in the main text, and reference Eq.~(...) where it is used.
%   • Use sub-section headings like
%       \section*{A.1 Derivation of ...} so the appendix reads as a
%       mini-textbook.
%
% Conventions:
%   • Use \appendix or follow the surrounding chapter numbering. Do NOT
%     mix numbered (\section) and unnumbered (\section*) indiscriminately.
%   • Equations inside appendices can use A.1, A.2, ... numbering. If you
%     prefer global Eq. labels, use \eqref{eq:appA:...} pattern.
%   • Cross-reference back into the main text where the formula is invoked.
% ==============================================================================

\chapter{Supplementary Derivations}
\label{app:derivations}

\section{Derivation of the Hartree--Fock Exchange Term}
In this appendix, we outline the mathematical details of the Hartree-Fock exchange energy term, which is utilized in hybrid exchange-correlation functionals such as HSE06. The exact exchange energy for a closed-shell system is given by:
\begin{equation}
    E_x^{\text{HF}} = -\frac{e^2}{2} \sum_{i,j} \iint \frac{\psi_i^*(\mathbf{r}_1) \psi_j(\mathbf{r}_1) \psi_j^*(\mathbf{r}_2) \psi_i(\mathbf{r}_2)}{4\pi\epsilon_0 |\mathbf{r}_1 - \mathbf{r}_2|} d\mathbf{r}_1 d\mathbf{r}_2
\end{equation}

In hybrid functionals, this term is split into short-range (SR) and long-range (LR) parts using an error function:
\begin{equation}
    \frac{1}{r} = \frac{\text{erfc}(\mu r)}{r} + \frac{\text{erf}(\mu r)}{r}
\end{equation}
where $\mu$ is the screening parameter. The short-range Hartree-Fock exchange is then mixed with DFT exchange to speed up the convergence in periodic systems.
